\section{Contour currents and angular transfer}\label{c7:contour}

The finite spin formula is regular at equal site coordinates, although its individual terms have poles there. We first replace the combined homogeneous value by a periodic Pfaffian contour expression. For each fixed number of sites, its limit as the period tends to infinity recovers the finite spin formula. We then identify the periodic expression, conditional on its Gaussian fields, with a fermionic trace propagating in the imaginary direction. Section~\ref{c7:boson} evaluates that trace, and Section~\ref{c7:stability} supplies the absolute estimate needed to integrate its reorganized series.

Use balanced groups of size $m$ each, with $m\in\mathbb N_0$. Write $\eta=\pi/4,\ p_0=2+\sqrt2,\ c=2-\sqrt2$, and put $z=e^{i\theta}$ and $h=z^{-2}$ for $\theta\in\mathbb C$. The logarithmic site coordinates are $x_j=\exp(\eta w_j)$. Normalize \eqref{j:finite-spin-sum} by its all-zero-spin term. Its value at homogeneous rapidities $w_j=0$ is
\[
P_0=p_0^{-2m}(4p_0^2)^{m^2}.
\]

\subsection{Defects and their pair kernels}

We split the factors between nonzero spins into a Schur Pfaffian and a Gaussian quadratic form. The accompanying one-site factor will provide the poles whose residues select the sites carrying nonzero spins.

First let both groups have the same generic distinct real coordinates $(w_i)_{i=1}^m$. Regard nonzero spin as a defect labeled $\delta=-s=\pm1$ with centered argument $v_i=w_i+i\delta/2$ (in shifts $i$ denotes the imaginary unit). Define
\[
\begin{split}
 V_-(w)&=\frac{e_1 e_7 e_3^2}{e_0 e_6 e_4^2},\qquad
 V_+(w)=\frac{e_1 e_7 e_5^2}{e_0 e_2 e_4^2}=\frac1{V_-(w+i)},\qquad e_k=1-q^k e^{\eta w},\\
 L(d)&=\frac{(y-x)^2(y^2+x^2)}{xy(y^2-\sqrt2xy+x^2)},\qquad
 J(d)=\frac{xy(y+x)^2}{(y^2+\sqrt2xy+x^2)^2},\qquad y/x=e^{\eta d}.
 \end{split}
\]
Then the normalized sum uses one factor
\[
z^\delta\,\frac{p_0\sqrt c}{4}\prod_{l\ne i}V_\delta(w_i-w_l)
\]
per defect and pair factors $L(v_j-v_i)^{\delta_i\delta_j}$ for the same group (species), $J(v_j-v_i)^{\delta_i\delta_j}$ across species, subject to equal total $\delta$.

Indeed remove powers $(xy)^{st}$ from same-group pairs of the table, $(xy)^{-st}$ across, and $x^{s^2}$ per site (total 1). For $s=+$ against two matched zeros the same-group ratio to baseline is $q^{-1}e_1e_7/(e_0e_6)$ and the cross ratio $q e_3^2/e_4^2$; for $s=-$ reflect $q$. At its own matched coordinate only the second contributes ($p_0/4$). Dividing the two-defect ratio by the two corresponding one-defect ratios gives $(\delta_i\delta_j)L(v_j-v_i)^{\delta_i\delta_j}$ or $(\delta_i\delta_j)J(v_j-v_i)^{\delta_i\delta_j}$ from the pair table (the argument shift multiplies $y/x$ by $q^{(s_i-s_j)/2}$). Total defect number is even, so the signs multiply to $(-1)^{n_-}$ with $n_-$ counting negative $\delta$. This cancels the site phase $i^{n_{\rm a}-n_{\rm b}}$ ($n_{\rm a},n_{\rm b}$ the defect counts per group) by equal total charges.

Introduce
\[
\begin{gathered}
a_*=\pi/6,\quad C(s)=\tanh(a_*s),\\
R^0(s)=C(s)^2/[C(s+i)C(s-i)],\\
G_d(s)=\log(R^0(s)/L(s)),\quad G_c(s)=-\log J(s).
\end{gathered}
\]
Take the real logarithms on the real line, with the limiting value at $s=0$, and then their analytic continuations. The matrix kernel $G$ has diagonal $G_d$ and off-diagonal $G_c$, with $G_d(0)=-\log A_*$, $A_*=3/(2c)$. At infinity $G_d(s)+\eta|s|,\ G_c(s)-\eta|s|$ decay exponentially. In Fourier notation $\widehat f(k)=\int e^{-ikv} f(v)\,dv$, the transforms of $G_d,G_c$ themselves at $k\ne0$, in units of $2\pi/k$, are
\[
g_d=\frac{\cosh4k+\cosh2k-\cosh3k}{\sinh4k}
 -\frac{\cosh3k-\cosh2k+\cosh k-1}{\sinh3k},\qquad g_c=\frac{1-2\cosh k}{\sinh4k}.
\]
The difference channel is used as covariance only with neutral tests. For instance $\log[2(\cosh(\eta v)-\cos(\eta b))]$, $0\le b\le4$, gives $-(2\pi/k)\cosh((4-b)k)/\sinh(4k)$ off zero, by expanding the log after subtracting $\eta|v|$ (sum of rational transforms is the elementary cosine series of $\cosh((4-b)k)$ on $0\le b\le4$). The same rule at period 6 gives the terms of $\log R^0$ using absolute values paired for conjugates. Constants at infinity fix transforms of integrable combinations; a residual linear $|v|$ part uses the usual neutral kernel by integrating $(1-\cos(kv))/k^2$. The covariance is positive in both channels; for $k>0,\ t=e^{-k}$,
\[
g_d+g_c=\frac{t^4(1-t)}{(1+t)(1+t^2)(1+t^4)(1-t+t^2)},\qquad g_c<0.
\]
Write $\widehat G_\pm=\widehat G_d\pm\widehat G_c$. The same formulas give
\begin{equation}\label{c7:gaussian-symbol-bounds}
\begin{aligned}
\widehat G_+(k)&=\pi/4+O(k^2),&
\widehat G_-(k)&=\pi/k^2+O(1) &&(k\longrightarrow0),\\
\widehat G_+(k)&=O(e^{-4|k|}/|k|),&
\widehat G_-(k)&=O(e^{-3|k|}/|k|) &&(|k|\longrightarrow\infty).
\end{aligned}
\end{equation}
Thus the common channel has the ordinary value $\widehat G_+(0)=\pi/4$, whereas the difference channel is applied to tests of total mass zero. The neutral formulas continue analytically for differences of arguments in $|\Im s|<3$.

For $k\ne0$, define multipliers
\[
E(k)=\frac{\cosh(3k/2)}{2\sinh(k/2)\cosh k},\qquad D(k)=\frac{k}{2\pi}E(k),
\]
and decaying functions $a,\lambda_0$ by odd real transforms, at $k>0$ given by
\[
\widehat a(k)=\frac{2\pi}{k}\frac{t^{1/2}(1-t)(1-t^4)}{(1+t^3)^2},\qquad
 \widehat\lambda_0=D\widehat a=\frac{\sinh k}{\cosh(3k/2)}.
\]
Both transforms have value zero at $k=0$. In particular $a$ and $\lambda_0$ are purely imaginary on the real line.
Define
\[
U_\delta(w)=V_\delta(w)\exp\!\left(\delta[(G_d+G_c)*\lambda_0](w+i\delta/2)\right).
\]
We claim that, for $|\Im v|<1/2$,
\[
U_-(v+i/2)=e^{a(v)},\qquad U_+(v-i/2)=e^{-a(v)}.
\]
For the first identity, note that $\log V_-(v+i/2)$ is decaying with zero winding near the real $v$ line: factors $1-e^{\eta(v+i b)}$ at $0<b<8$ have phase change $\eta b-\pi$, and the signed multiplicities at $1.5,7.5,3.5,0.5,6.5,4.5$ are respectively $(1,1,2,-1,-1,-2)$, balanced also in the sums of $b$. In fact if a sum of such logs (plus a constant) has multiplicities $n_b$ with $\sum n_b=0$, its derivative decays exponentially at both horizontal infinities and is $8i$-periodic with residues $n_b$ at $-ib$ modulo period. Shifting the derivative transform down by $8i$ (positive $k$) thus gives the log transform $-(2\pi/k)\sum_b n_b t^b/(1-t^8)$ there. This off-zero calculation also applies for bounded logs with step limits by taking the transform via differentiation. For our decaying log it gives the expression $(2\pi/k)t^{1/2}(1-t)(1-t^3)^2/(1-t^8)$ (extended oddly). This equals $(1+D(\widehat G_d+\widehat G_c))\widehat a$. The Fourier bounds give the stated strip continuation, and $V_+(w)=V_-(w+i)^{-1}$ gives the second identity. In particular
$U_+(w)=U_-(w+i)^{-1}$ meromorphically. The decays and strip continuations also follow by shifting Fourier integrals; the convolution in $U$ continues well beyond $|\Im|<2$.

\subsection{A fixed-site periodic Pfaffian}

Fix a horizontal period $T>0$. For a decaying function, subscript $T$ denotes the sum over translates by $T$; for a meromorphic function tending to $1$ exponentially at both horizontal infinities, it denotes the product over those translates. These sums and products converge locally uniformly on the strips used below, away from their stated poles. An integral $\int_T dx$ is over one full period and is unnormalized. For $k=2\pi n/T$, the coefficient of $e^{ikx}$ in a periodized decaying function is $\widehat f(k)/T$.

Let $P_\pm$ be the orthogonal projections on the spans of $(1,1)^t$ and $(1,-1)^t$. Take centered real Gaussian oscillators
\[
\chi_\sigma(v)=\sum_{k=2\pi n/T\ne0}\chi_{\sigma,k}e^{ikv},
\qquad \chi_{\sigma,-k}=\overline{\chi_{\sigma,k}},
\]
whose covariance between the coefficients at $k,-k$ is $\widehat G(k)/T$, and whose other mode covariances vanish. Positivity follows from the two channels above. By \eqref{c7:gaussian-symbol-bounds},
\[
\mathbb E\sum_{k\ne0}|\chi_{\sigma,k}|e^{\rho|k|}<\infty
\qquad(0\le\rho<3/2).
\]
Consequently these series are almost surely analytic on $|\Im v|<3/2$, with absolute convergence on every smaller closed strip.

Choose a centered common Gaussian vector with covariance $\widehat G_+(0)P_+/T$ and a uniform vector $(t_0,-t_0)$ with $0\le t_0<2\pi$, independently of each other and of $\chi$. Their sum is denoted by $(\varphi_{01},\varphi_{02})$. Independently of all preceding variables, choose
\[
\boldsymbol\ell=(\ell_1,\ell_2)=(j,-j),\quad j\in8\mathbb Z,\qquad
\mathbb P(j)=D_\ell^{-1}e^{-4\pi j^2/T},\quad
D_\ell=\sum_{j\in8\mathbb Z}e^{-4\pi j^2/T},
\]
and set
\[
\varphi_\sigma(v)=\varphi_{0\sigma}+(2\pi/T)\ell_\sigma v+\chi_\sigma(v).
\]
We write $\varphi(v)=(\varphi_1(v),\varphi_2(v))^t$.

For each species use a Wick expectation: the expectation of an ordered list of an even number of fields is the Pfaffian of their two-point entries, and odd expectations are zero. Its two-point function is
\[
\langle\psi(u)\psi(v)\rangle_T=C^{[T]}(v-u),\qquad
C^{[T]}(d)=\sum_{n\in\mathbb Z}
 \frac{\pi}{a_*T\sin(\pi(d+3i+6in)/T)}.
\]
For fixed $T$ the image series converges locally uniformly away from its poles. It is odd, $6i$-periodic, horizontally antiperiodic, and has residue $1/a_*$ at $3i$.

To define the contour expression also at coalescing sites, choose
\[
0<r<\min\{1/8,T/8\},\qquad |w_l|<r/4\quad(1\le l\le m),
\]
initially with distinct real $w_l$, and let $\gamma_r$ be the positive circle $|w|=r$. Put $W=(w_1,\ldots,w_m)$ and
\[
\Lambda_W(x)=\sum_{l=1}^m\lambda_{0,T}(x-w_l).
\]
Define, at first as its Wick expansion,
\begin{equation}\label{c7:contour-current}
\mathcal C_{\sigma,T}(W)=
\left\langle\exp\left(
a_*\sum_{\delta=\pm1}\oint_{\gamma_r}\frac{dw}{2\pi i}
 (-z^\delta)e^{i\delta\varphi_\sigma(w+i\delta/2)}
 \prod_{l=1}^m U_{\delta,T}(w-w_l)\,
 \psi(w+3i)\psi(w+i\delta)\right)\right\rangle_T .
\end{equation}
The same list $W$ is used for both species. The normalized expression is
\begin{equation}\label{c7:regulated-contour}
\mathfrak C_T(W;\theta)=
\frac{\mathbb E\!\left[
 e^{i\int_T\Lambda_W(x)(\chi_1+\chi_2)(x)\,dx}
 \mathcal C_{1,T}(W)\mathcal C_{2,T}(W)\right]}
{\mathbb E e^{i\int_T\Lambda_W(x)(\chi_1+\chi_2)(x)\,dx}}.
\end{equation}
At $W=0^m=(0,\ldots,0)$ write $\mathcal C_\sigma=\mathcal C_{\sigma,T}(0^m)$ and $\mathfrak C_{m,T}=\mathfrak C_T(0^m;\theta)$.
For comparison with the finite formula put
\[
x_l=e^{\eta w_l},\qquad
X(W)=(x_1,\ldots,x_m,x_1,\ldots,x_m),\qquad
P_0(W)=\prod_{i<j}p_{00}(x_i,x_j)^2
       \prod_{i,j}p'_{00}(x_i,x_j).
\]
This is the all-zero term for the matched list. The displayed domain avoids its zeros and poles, and its value at $W=0^m$ is $P_0$.
The first and second copies in $X(W)$ belong to $A$ and $B$, respectively.

\begin{proposition}[Fixed-site regulated Pfaffian]\label{c7:regulated-pfaffian}
For every fixed $T>0$, $m\in\mathbb N_0$, and $\theta\in\mathbb C$, on the full-probability event of analytic samples described above, the Wick expansion in \eqref{c7:contour-current} terminates after order $m$ in each species. Its value is independent of the admissible circle enclosing $W$, and its contour coefficients extend holomorphically to the coalescing lists $|w_l|<r/4$. The average \eqref{c7:regulated-contour} is holomorphic there and its denominator is nonzero. In particular $\mathfrak C_{0,T}=1$.

For fixed $m\ge1$ and $\theta\in\mathbb C$,
\begin{equation}\label{c7:regulated-limit}
\lim_{T\longrightarrow\infty}\mathfrak C_{m,T}
=\frac{N(1,\ldots,1)^2}{P_0}\,
G_{A,B}(h+h^{-1}),\qquad |A|=|B|=m.
\end{equation}
This limit is locally uniform in $\theta$.
\end{proposition}

\begin{proof}
For any distinct complex list with $|w_l|<r/4$, the only poles of each current variable inside $\gamma_r$ are the simple site poles. The Pfaffian is analytic there in all current variables: its nearest inter-current poles require an imaginary difference $\pm1$ between the unshifted variables, whereas their differences in this disk have imaginary part less than $2r<1$. If two residues in the same species select the same site, their first fields $\psi(w+3i)$ coincide, so the Pfaffian vanishes. Thus at most $m$ currents occur in each species. Cauchy's theorem applied to each coefficient also shows independence of the admissible circle. On $\gamma_r$ every coefficient is holomorphic in the $w_l$ for $|w_l|<r/4$. The coefficients of higher orders vanish there by continuity from distinct lists, proving termination at confluence as well. If $m=0$, there are no enclosed site poles and $\mathcal C_\sigma=1$.

We next compute the Gaussian factor of any one of these finitely many terms. If its current variables are $\omega_a\in\gamma_r$, let $v_a=\omega_a+i\delta_a/2$ and set
\[
q_a=\delta_a e_{\sigma_a},\quad
Q(k)=\sum_aq_a e^{-ikv_a},\quad Q_0=Q(0),\quad
V_Q=\sum_aq_a v_a,\quad \mathbf1=(1,1)^t,\quad \mathbf e_-=(1,-1)^t,
\]
where $e_1,e_2$ are the species basis vectors. The Fourier transform of the source is
\[
S_W(k)=\widehat\lambda_0(k)\sum_l e^{-ikw_l}\mathbf1.
\]
For $k=2\pi n/T\ne0$, the exact ratio, including the constant and winding modes, is
\begin{equation}\label{c7:gaussian-coefficient}
\begin{aligned}
&\frac{\mathbb E\exp\!\left(
 i\int_T\Lambda_W(\chi_1+\chi_2)+i\sum_aq_a^t\varphi(v_a)\right)}
 {\mathbb E\exp\!\left(i\int_T\Lambda_W(\chi_1+\chi_2)\right)}
\\
&\quad=
\exp\!\left\{-\frac1{2T}\sum_{k\ne0}
 \left[Q(-k)^t\widehat G(k)Q(k)
       +2S_W(-k)^t\widehat G(k)Q(k)\right]\right\}
\\
&\qquad\quad{}\times
\mathbf1_{\{(Q_0)_1=(Q_0)_2\}}\,
\exp\!\left[-\frac{\widehat G_+(0)}{2T}Q_0^tP_+Q_0\right]
\frac1{D_\ell}\sum_{j\in8\mathbb Z}
 e^{-4\pi j^2/T+(2\pi i j/T)\mathbf e_-^tV_Q}.
\end{aligned}
\end{equation}
All pairings in this formula are bilinear, without complex conjugation. The uniform mode gives the charge constraint exactly at each $T$. The oscillator part is the elementary centered Gaussian exponential identity, first on finitely many modes and then in the limit. Its denominator is
\[
\exp\!\left[-\frac1{2T}\sum_{k\ne0}
 S_W(-k)^t\widehat G(k)S_W(k)\right],
\]
which is nonzero; for real $W$ it is positive. The covariance sums converge absolutely, since $|\Im v_a|\le1/2+r<3/2$, while $\widehat\lambda_0(k)=O(e^{-|k|/2})$ and $|\Im w_l|<r/4$ control the source. Together with \eqref{c7:gaussian-symbol-bounds}, these estimates give an exponential margin and finite exponential moments uniformly on the compact contours for each fixed $T$. They also make the coefficient sums holomorphic in $W$. The winding sum is dominated by $e^{-4\pi j^2/T+C|j|/T}$. Thus the Gaussian and contour integrations in the finite expansion are justified.

The cross term with $S_W$ in \eqref{c7:gaussian-coefficient} is
\[
-\sum_a\delta_a\sum_l
 [(G_d+G_c)*\lambda_0]_T(v_a-w_l).
\]
The zero Fourier coefficient of $\lambda_0$ vanishes, so the periodized convolution here is exactly the periodic convolution arising from the source. The definition of $U_\delta$ and the normally convergent periodized products therefore show that this cross term converts every $U_{\delta,T}$ factor to $V_{\delta,T}$ already at finite $T$.

It remains to take the limit of the combined quadratic in \eqref{c7:gaussian-coefficient}. On the terms surviving the uniform mode, $P_-Q(0)=0$, and hence $P_-Q(k)=O(k)$ uniformly on the contours. This cancels the $k^{-2}$ singularity of $\widehat G_-$ inside the complete quadratic. The common channel is regular at zero. At high frequency the pairwise imaginary differences have absolute value at most $1+2r<3$. Consequently the Riemann sums converge uniformly to the neutral Fourier integral, giving
\begin{equation}\label{c7:neutral-limit}
\frac1T\sum_{k\ne0}Q(-k)^t\widehat G(k)Q(k)
\longrightarrow
\sum_{a,b}q_a^tG(v_b-v_a)q_b.
\end{equation}
Here the common channel uses its ordinary Fourier inverse, and the difference channel uses its zero-total-mass interpretation described above. Combining the quadratic before the limit is essential: the periodic difference-channel variance at a single point grows as $\pi T/12$. Formula \eqref{c7:neutral-limit}, rather than separate limits of its discrete self and pair terms, permits the self and pair decomposition of the limiting $G$ form.

The constant Gaussian factor in \eqref{c7:gaussian-coefficient} tends to $1$. The winding factor does as well, uniformly for bounded $V_Q$: its distribution has $\mathbb E|j|=O(\sqrt T)$ and each fixed exponential moment of $|j|/\sqrt T$ is bounded as $T$ grows. The periodized scalar factors tend locally uniformly to their unperiodized versions. For the fermion kernel, pair the image terms with imaginary shifts $\pm(3+6n)i$. The sine derivative bound gives a uniform $O((n+1)^{-2})$ bound on each paired term for $d$ in a compact set off the poles. Termwise limits and the elementary partial fractions of $\tanh$ give $C^{[T]}(d)\to C(d)$ there.

The limiting Wick Pfaffian has the Schur product form
\[
\left\langle\psi(u_1)\cdots\psi(u_{2s})\right\rangle
=\prod_{a<b}C(u_b-u_a).
\]
Indeed, after $y_a=e^{2a_*u_a}$, its kernel is $(y_b-y_a)/(y_b+y_a)$. Alternating Vandermonde divisibility and induction at a pairing pole prove that its Pfaffian is the full pair product. Within a current the product gives $C(i\delta-3i)=1/C(i\delta)$. Between two currents it gives $R^0(v_b-v_a)^{\delta_a\delta_b}$. In the representation \eqref{c7:neutral-limit}, each self term gives $\exp(-G_d(0)/2)=\sqrt{A_*}$, and the pair terms change these factors to $L^{\delta_a\delta_b}$ within a species and $J^{\delta_a\delta_b}$ across species. Finally
\[
C(i\delta)=i\delta/\sqrt3,\qquad
-a_*\operatorname{Res}_{w=0}V_\delta(w)
 \frac{\sqrt{A_*}}{C(i\delta)}
=\frac1{2\sqrt c}=\frac{p_0\sqrt c}{4},
\]
using $\operatorname{Res}_0V_\delta=-i\delta\sqrt2/(4\eta)$ and $cp_0=2$. For distinct sites these are exactly the one-defect and pair factors derived from the finite spin sum.

For the confluence to the homogeneous value, keep one fixed circle $\gamma_r$ with $0<r<1/8$, let $|w_l|\le r/8$, and take $T>8r$ to infinity. All the covariance, image, and scalar convergence just used is uniform on this circle and polydisk. These local estimates permit the limiting contour expression to be evaluated at $W=0^m$. At distinct real $W$ the residue calculation identifies it with the normalized finite spin expression. By Section~\ref{c7:finite}, the numerator $N(X(W))^2\mathscr G(X(W);A,B)$ has a polynomial continuation and $P_0(W)$ is nonzero here. Holomorphic uniqueness in the auxiliary coordinates therefore identifies the value at $W=0^m$, giving \eqref{c7:regulated-limit}. Compact sets of $\theta$ cause only bounded factors in this finite expansion, proving the stated local uniformity.
\end{proof}

The proposition takes the infinite-period limit at fixed $m$, before introducing any infinite expansion of trace operators. The identities in the rest of this section hold for each fixed $T>0$. Their use as an absolutely integrated horizontal gas will require the later stability estimate, which is needed only for sufficiently large integer block periods.

\subsection{A conditional angular trace}

We now fix $T>0$, $m$, $\theta$, and one analytic sample of the fields, and use homogeneous sites. The contraction rule is the normal-ordering principle of Wick~\cite{Wick1950}; here we compute its periodic covariance and trace realization. Suppress the species index. Fold the imaginary coordinate into six sheets, indexed cyclically by $j=1,\ldots,6$, and write $j^*=j+3$. On the fermionic exterior space take
\[
\{\Psi_{j,r},\Psi_{l,s}\}=\delta_{j^*,l}\delta_{r,-s},
\qquad
\Psi_{j,r}^*=\Psi_{j^*,-r},
\qquad r,s\in\tfrac12+\mathbb Z,
\qquad
\Psi_j(v)=\sum_r\Psi_{j,r}y^{-r},\quad y=e^{2\pi i v/T}.
\]
The modes with $r>0$ annihilate the vacuum. Let $L_0$ be total creation energy, let $\mathbf q=e^{-2\pi/T}$, and let $\mathcal U$ fix the vacuum and send $\Psi_j$ to $-\Psi_{j-1}$ under conjugation. Put $\Gamma=\mathbf q^{L_0}\mathcal U$.

For the off-diagonal roots below define
\[
X_{bc}(f)=-\sum_{n\in\mathbb Z}f_n
 \sum_{r+s=n}\Psi_{b^*,r}\Psi_{c,s}
=-\int_T\frac{dv}{T}\,f(v)\Psi_{b^*}(v)\Psi_c(v),
\qquad
J_{bc}=\mathrm e_{bc}-\mathrm e_{c^*b^*},
\]
where $f(v)=\sum_nf_ny^n$ is analytic and periodic and the integral means its modewise Fourier integral. We only use $b\notin\{c,c^*\}$. The notation $X_{bc}(f;\gamma)$ uses a shifted horizontal line $\gamma$; it gives the same Fourier integral when $f$ is analytic between the lines. As a formal mode series, the commutator of $X_{bc}(f)$ with the column of fields is $fJ_{bc}$ times that column. Thus the matrices of conjugations multiply in the reverse order from the operators.

For the fixed sample set
\[
b(v)=z^{-1}e^{-i\varphi(v)}[U_{-,T}(v+i/2)]^m,
\qquad b_\pm(v)=b(v\pm i).
\]
This function is periodic, since $\ell\in8\mathbb Z$. For $U_{-,T}(v+i/2)$, the nearest zero and pole occur at $v=i/2$ and $v=-i/2$ modulo $T$, respectively; the next divisors occur at $v=\pm3i/2$ modulo $T$. Its exponential factor is nonzero wherever it is analytic. Combining these facts with the sample strip $|\Im v|<3/2$ shows that $b,b_+,b_-$ and their inverses are analytic on $|\Im v|<1/2$, bounded on each smaller closed strip over one period. These are conditional bounds depending on the sample and the fixed parameters.

\begin{lemma}[Conditional angular trace]\label{c7:angular-trace}
Fix $T>0$, $m\in\mathbb N_0$, $\theta\in\mathbb C$, and a homogeneous analytic sample as above. The normalizing trace is
\[
Z_\Gamma:=\operatorname{Tr}\Gamma
=\prod_{r\in\frac12+\mathbb N_0}(1-\mathbf q^{6r})>0.
\]
For $s\ge1$, if the folded coordinates have strictly increasing imaginary parts with total range less than $1$, then
\begin{equation}\label{c7:angular-wick}
\left(-\frac{2\pi i}{a_*T}\right)^s
\frac{\operatorname{Tr}\bigl(\Gamma\Psi_{j_1}(v_1)\cdots
 \Psi_{j_{2s}}(v_{2s})\bigr)}{Z_\Gamma}
=
\operatorname{Pf}_{1\le a,b\le2s}
 C^{[T]}\bigl(v_b-v_a+i(j_b-j_a)\bigr).
\end{equation}
The trace in this display denotes its absolutely convergent ordered mode sum. For $s=0$ both sides are $1$. The formula at other nonsingular coordinates is its meromorphic continuation; odd products vanish.

Let $P_N$ project onto the energy-$N$ subspace, for $N\in\frac12\mathbb N_0$. For a fixed finite word $Y$ of exponentials of the currents $X_{bc}(f)$, with periodic coefficients analytic on a common strip of positive width, there are $C,c>0$ such that
\begin{equation}\label{c7:graded-current-bound}
\|P_{N'}YP_N\|\le C e^{C\sqrt{N+1}}e^{-c|N'-N|}.
\end{equation}
The constants can be taken locally uniform in numerical parameters multiplying the currents in the exponentials. In particular the energy-level trace with $\Gamma$ is absolutely convergent, as are its exponential series of scalar trace coefficients.

For the homogeneous contour expression \eqref{c7:contour-current} and the fixed analytic sample,
\begin{equation}\label{c7:four-current-trace}
\mathcal C_\sigma
=\frac1{Z_\Gamma}\operatorname{Tr}_{\rm gr}\!\left(
 \Gamma e^{X_{23}(b_+^{-1})}e^{X_{21}(b)}
 e^{-X_{12}(b^{-1})}e^{-X_{16}(b_-)}\right).
\end{equation}
Here the absolutely convergent trace furnished by \eqref{c7:graded-current-bound} is
\[
\operatorname{Tr}_{\rm gr}(\Gamma Y)
=\sum_N\mathbf q^N
 \operatorname{tr}_{P_N\mathcal F}(\mathcal U P_NYP_N),
\]
where $\mathcal F$ is the fermionic exterior space. The convergence assertions concern this fixed sample.
\end{lemma}

\begin{proof}
The operator $\mathbf q^{L_0}$ is trace class because
$\prod_{r>0}(1+\mathbf q^r)^6<\infty$, and $\mathcal U$ is unitary. The trace of the exterior space at each positive energy $r$ is the determinant of $I$ plus the one-particle action. That action is $-\mathbf q^r$ times a six-cycle, so its determinant is $1-\mathbf q^{6r}$. Since $\sum_{r>0}\mathbf q^{6r}<\infty$, this proves the product for $Z_\Gamma$ and its positivity.

Cyclicity for bounded individual CAR modes gives, for $r>0$,
\[
\frac{\operatorname{Tr}(\Gamma\Psi_{j,r}\Psi_{l,-r})}{Z_\Gamma}
=\sum_{n\ge0}\mathbf q^{rn}\delta_{l-n,j^*}.
\]
Indeed the left side, denoted $C_{jl}$, obeys
$C_{jl}=\delta_{l,j^*}+\mathbf q^r C_{j+1,l}$.
The CAR then give
$-\sum_{n\ge1}\mathbf q^{-rn}\delta_{l+n,j^*}$ for $r<0$.
Other energy pairings vanish. For $0<\Im(v-u)<1$, summing the half-integer modes yields
\[
\frac{\operatorname{Tr}(\Gamma\Psi_j(u)\Psi_l(v))}{Z_\Gamma}
=\frac{i}{2}
 \sum_{\substack{n\in\mathbb Z\\n\equiv l-j-3\pmod6}}
 \frac1{\sin(\pi(v-u+in)/T)}.
\]
Here the positive and negative geometric sums converge respectively on the two sides of the annulus. Multiplication by $-2\pi i/(a_*T)$ gives
$C^{[T]}(v-u+i(l-j))$. Moving one mode cyclically in a higher product, contracting it by the CAR, and solving the same six-step recurrence gives Wick's Pfaffian rule. For two ordered heights $\eta_a<\eta_b$, direct contractions have decay from $\eta_b-\eta_a$, and contractions wrapping through $\Gamma$ have decay from $1-(\eta_b-\eta_a)$. Both are positive in the stated strip, which makes the finite products of mode sums absolute. This proves \eqref{c7:angular-wick}; graded alternation gives its meromorphic continuation.

We give the estimates needed for the integrated currents. Suppose $f$ is bounded and analytic on a closed strip of width $\rho>0$. Its coefficients obey
$|f_n|\le C e^{-2\pi\rho|n|/T}$. The quadratic mode with $r+s=n$ changes energy by $-n$. On energy $N$ its norm is at most
\[
C(\sqrt{N+1}+|n|).
\]
For the creator-annihilator part this is the second-quantized norm of a bounded one-particle shift times the particle count, which is $O(\sqrt{N+1})$ by exclusion. The two-creator and two-annihilator parts have $O(|n|)$ terms. Normal ordering adds no contraction for these roots.

We also need a bound before summing the constituent mode pairs. On an occupation state of energy $N$, at most $O(N+|n|+1)$ pairs $r+s=n$ can act: in the mixed-sign part the annihilator must meet an occupied mode, and the same-sign part has $O(|n|)$ pairs. Each constituent has norm at most one. This coarser count, the Fourier exponential, and the energy damping below give absolute sums of occupation-basis terms for any fixed monomial of currents.

For a product of $s$ quadratic modes, let $u$ be the sum of the absolute energy movements. Absorbing the individual $|n|$ factors into part of the Fourier exponential leaves a bound
$C^s(1+\sqrt N+\sqrt u)^s$. Absorbing another part and summing the movements gives
\[
C_1^s(1+\sqrt N+\sqrt s)^s e^{-c|N'-N|}
\]
before the exponential-series factorials. Summing in $s$ gives \eqref{c7:graded-current-bound}: splitting $(1+\sqrt N+\sqrt s)^s$ into the two terms with bases $1+\sqrt N$ and $\sqrt s$, and using the factorial, bounds the sum by $e^{C\sqrt{N+1}}$. Finite-energy dimensions are bounded by $e^{C\sqrt{N+1}}$, as follows by evaluating the exterior generating product at nome $e^{-1/\sqrt{N+1}}$. The factor $\mathbf q^N$ therefore makes the energy-level trace absolute. The same estimates are locally uniform in numerical exponential parameters, so commutators, differentiation in those parameters, and the conjugations below can be verified by convergent coefficients. The argument also applies to the nonzero normally ordered Cartan modes used below.

It remains to relate the contour to those currents. Use $0<\varepsilon<1/8$ and expand $\gamma_r$ to the boundary of the band $|\Im w|<1/2-\varepsilon$ in $\mathbb C/T\mathbb Z$. No inter-current pole is crossed, since all unshifted differences have imaginary part strictly between $-1$ and $1$. The Gaussian arguments remain inside $|\Im v|<3/2$. Each bilinear is periodic horizontally, so the vertical sides cancel. On the top set $v=w-i/2$ and on the bottom set $v=w+i/2$. The identification in \eqref{c7:angular-wick} is
\[
\psi(v+i(j-3/2))
\longleftrightarrow
\left(-\frac{2\pi i}{a_*T}\right)^{1/2}\Psi_j(v).
\]
A positive boundary runs west along the top and east along the bottom. Using $U_+(w)=U_-(w+i)^{-1}$, the four contributions are
\begin{center}
\small
\begin{tabular}{@{}llll@{}}
\toprule
boundary and defect&folded line&ordered sheets&current\\
\midrule
top, $\delta=-1$&$w=v+i/2,\ \Im v=-\varepsilon$&$(5,1)$&$X_{21}(b)$\\
top, $\delta=+1$&$w=v+i/2,\ \Im v=-\varepsilon$&$(5,3)$&$X_{23}(b_+^{-1})$\\
bottom, $\delta=+1$&$w=v-i/2,\ \Im v=\varepsilon$&$(4,2)$&$-X_{12}(b^{-1})$\\
bottom, $\delta=-1$&$w=v-i/2,\ \Im v=\varepsilon$&$(4,6)$&$-X_{16}(b_-)$\\
\bottomrule
\end{tabular}
\end{center}
For example, the coefficient on the unoriented top line for $\delta=-1$ is $-b(v)$, and reversing that line changes it to $b(v)$. The pair of physical field multipliers changes $a_*\,dv/(2\pi i)$ to $-dv/T$, exactly the normalization in $X_{bc}$.

Within either line there are no dual field pairs, so its bilinears have no internal contraction poles and may be separated infinitesimally. The top folded line is below the bottom folded line and their range is $2\varepsilon<1$, as required by \eqref{c7:angular-wick}. Match first each coefficient after inserting one numerical parameter multiplying all contour currents. The contour deformation and Wick rule give the corresponding coefficient of the four current exponentials; exchanging entire bilinears costs no sign because they are even. The two top roots commute, and so do the two bottom roots. The graded estimate makes the trace series entire in the numerical parameter. Summing the matched coefficients at parameter $1$ proves \eqref{c7:four-current-trace}.
\end{proof}

\subsection{The seam operator and outgoing coefficients}

The four-current trace can be reorganized so that one operator carries the exchange of sheets $1,2$. Define
\[
W_b=e^{-X_{12}(b^{-1})}e^{X_{21}(b)}e^{-X_{12}(b^{-1})}.
\]
It acts on the formal field series by the matrix $P_b$ whose blocks on sheets $1,2$ and $4,5$ are
\[
P_b|_{1,2}=\begin{pmatrix}0&-b^{-1}\\ b&0\end{pmatrix},
\qquad
P_b|_{4,5}=\begin{pmatrix}0&-b\\ b^{-1}&0\end{pmatrix}
=(P_b|_{1,2})^{-t},
\]
and which fixes sheets $3,6$. The exact operator identity is
\begin{equation}\label{c7:seam-factorization}
\begin{aligned}
&e^{X_{23}(b_+^{-1})}e^{X_{21}(b)}
 e^{-X_{12}(b^{-1})}e^{-X_{16}(b_-)}
\\
&\qquad=
e^{X_{26}(bb_-)}e^{X_D}W_b
 e^{-X_{13}((bb_+)^{-1})},\\
&X_D=X_{12}(b^{-1})+X_{13}((bb_+)^{-1})
       -X_{16}(b_-)+X_{15}(b_-/b_+).
\end{aligned}
\end{equation}
We verify the scalar as well as the field action. Put
$A_{\rm r}=X_{23}(b_+^{-1})$, $B_{\rm r}=X_{26}(bb_-)$, and
$C_{\rm r}=X_{12}(b^{-1})$. Isolating the defining three-root word for $W_b$ and conjugating the left field of $X_{16}$ by $W_b$ writes the left side as
$e^{A_{\rm r}}e^{C_{\rm r}}e^{B_{\rm r}}W_b$.
The roots $A_{\rm r},B_{\rm r}$ commute, with no central term: their only possible single contraction leaves two copies of $\Psi_5$, and two contractions are impossible. Thus the middle exponent after moving them past it is
$e^{-\operatorname{ad}B_{\rm r}}e^{\operatorname{ad}A_{\rm r}}C_{\rm r}$.
In this conjugation $\Psi_4$ stays fixed and its partner changes by
\[
\Psi_2\longmapsto
\Psi_2+b_+^{-1}\Psi_3-bb_-\Psi_6+(bb_-/b_+)\Psi_5.
\]
Multiplying by $b^{-1}$ gives exactly $X_D$. Also
$W_b^{-1}A_{\rm r}W_b=-X_{13}((bb_+)^{-1})$.
Every transformed partner in these steps is nondual to its first field, so normal ordering produces no scalar contraction. The graded convergence in Lemma~\ref{c7:angular-trace} makes these conjugations identities of the operators used in the trace. The field calculation therefore fixes the scalar as well as the matrices $J_{bc}$ and proves \eqref{c7:seam-factorization}.

The part of $W_b$ depending on the oscillators will be evaluated in Section~\ref{c7:boson}. Periodizing $U_-(v+i/2)=e^{a(v)}$ gives $U_{-,T}(v+i/2)=e^{a_T(v)}$ on $|\Im v|<1/2$. Hence write
\[
b=b_w e^{\mathsf A},\qquad
b_w=b_0y^{-\ell},\qquad b_0=e^{-i(\theta+\varphi_0)},
\qquad \mathsf A=ma_T-i\chi=\sum_n\mathsf A_ny^n .
\]
Here $\mathsf A_0=0$ because $\widehat a(0)=0$ and $\chi$ has zero spatial mean. For $j=1,2,3$ use the normally ordered currents
\[
h_{j,n}=\sum_{r+s=n}:\Psi_{j,r}\Psi_{j^*,s}:,
\qquad [h_{j,n},\Psi_j(v)]=y^n\Psi_j(v).
\]
Then
\begin{equation}\label{c7:cartan-seam}
W_b=W_{b_w}
 \exp\!\left(\sum_n\mathsf A_n(-h_{1,n}+h_{2,n})\right).
\end{equation}
Indeed, with $K=\frac12\sum_n\mathsf A_n(h_{1,n}-h_{2,n})$, conjugating the defining root word for $W_{b_w}$ by $e^K$ changes its coefficients to those of $W_b$. The action of $W_{b_w}$ interchanges the nonzero $h_{1,n},h_{2,n}$, so $W_{b_w}KW_{b_w}^{-1}=-K$ and
$e^KW_{b_w}e^{-K}=W_{b_w}e^{-2K}$.
Only nonzero modes enter. At zero mode the Cartan difference would instead transform to
$-(h_{1,0}-h_{2,0})+2\ell$, which explains the use of $\mathsf A_0=0$.

Finally we put each fixed coefficient of the remaining three exponentials on the outgoing side of $W_b$. On the finite-energy core the mode identity is
\[
\Gamma\Psi_{j,r}\Gamma^{-1}=-\mathbf q^{-r}\Psi_{j-1,r}.
\]
Equivalently, as a formal mode series,
\[
\Gamma\Psi_j(v)\Gamma^{-1}=-\Psi_{j-1}(v+i).
\]
The corresponding current rule, including its scalar argument, is
\begin{equation}\label{c7:cyclic-current}
\Gamma X_{bc}(f;\gamma)\Gamma^{-1}
=X_{b-1,c-1}\bigl(f(\,\cdot-i\,);\gamma+i\bigr).
\end{equation}
The two field signs cancel and the orientation is unchanged. In a trace this identity is first applied to fixed modes; the movement estimates in the proof of Lemma~\ref{c7:angular-trace}, before factorial summation, give absolute sums for each fixed monomial and permit their reindexing and Fourier integration.

More explicitly, let $Y_L,Y_D,Y_U$ be fixed monomials from the leftmost, middle, and rightmost exponentials of \eqref{c7:seam-factorization}. Cyclicity gives
\[
\operatorname{Tr}_{\rm gr}(\Gamma Y_LY_DW_bY_U)
=\operatorname{Tr}_{\rm gr}
 \bigl(\Gamma W_bY_U(\Gamma Y_L\Gamma^{-1})
                       (\Gamma Y_D\Gamma^{-1})\bigr).
\]
The denominator remains $Z_\Gamma$. Call the rightmost root $U$, the cycled leftmost root $L$, and retain $D$ for the middle sum. Their outgoing order is $U,L,D$. One may initially put the unshifted $U$ pairs just above folded height $0$, and the incoming $L,D$ pairs on lines $-\rho_L,-\rho_D$ with $0<\rho_D<\rho_L<1/2$. After \eqref{c7:cyclic-current} these latter lines have heights $1-\rho_L,1-\rho_D$. Further infinitesimal separations realize strict order
\[
0<\eta_U<\eta_L<\eta_D<1
\]
for each fixed coefficient, within the original scalar analytic strips. Section~\ref{c7:boson} evaluates these ordered coefficients and establishes its further contour displacements from their paired analytic expressions.

Thus the Gaussian expectation of the total trace is defined by its equality to the finite contour polynomial. Passing that expectation through the reorganized infinite expansion requires the absolute estimate of Section~\ref{c7:stability}. The conditional trace identities here hold for every fixed $T>0$; the later integrated gas identity uses sufficiently large integer periods to which that estimate applies.
